\section{Dyadic cokernel bounds and a common exceptional window}
\label{sec:sr-dyadic-bounds}
The generator and descent results of
Sections~\ref{sec:sr-dyadic-generators}--\ref{sec:sr-dyadic-descent}
permit a quantitative argument with fixed constants.
We prove that argument here. Its last step is conditional on a
single high-twist cohomology bound, stated explicitly in
Definition~\ref{def:sr-K2}; that bound is not proved by the
linear algebra below.

Fix the principal primes \(q,d\), the field
\(L=\mathbb Q(i,\sqrt{-q},\sqrt{-d})\), its cyclotomic extension
\(L_\infty\), and the split unramified module \(M=X_2\).
We have a continuous surjection \(\mathbb Z_2^{24}\to M\).
The generator \(5\) of the cyclotomic group induces the continuous
automorphism \(\tau\) of \(M\).
For \(a\geq1,m\geq0\), write
\[
 c_m(a)=\operatorname{length}_{\mathbb Z_2}
           \bigl(M/((5^a\tau)^{2^m}-1)M\bigr).
\]
Lengths may be infinite until finiteness has been established.
Proposition~\ref{prop:sr-dyadic-descent} identifies this with
\(s_{2,L_m}(a)\), and \eqref{eq:sr-dyadic-return} bounds
\(s_{2,E_h}(a)\) by \(1+c_h(a)\).
In this section the fixed depth is \(j_0=7\).
It is distinct from the exponent \(J\) controlling the analytic window.

\subsection{A free lift and its uniform congruence}
Put \(r=\dim_{\mathbb F_2}(M/2M)\leq24\).
Lifts of an \(\mathbb F_2\)-basis generate \(M\) over
\(\mathbb Z_2\). Indeed their closed span is the compact image of
\(\mathbb Z_2^r\); the quotient \(Q\) satisfies \(Q=2Q\).
Then \(Q=2^nQ\) for every \(n\), while
\(\bigcap_n2^nQ=0\) by the finite-quotient argument in
Lemma~\ref{lem:sr-dyadic-bidual}. Thus \(Q=0\).
If \(r=0\), \(M=0\) and the assertions below are immediate.
Assume henceforth \(r>0\), and choose a surjection
\(\pi:\mathbb Z_2^r\to M\) lifting such a basis.

Choose a matrix \(A\in M_r(\mathbb Z_2)\) such that
\(\pi A=\tau\pi\), by lifting the images of the basis vectors.
Modulo two, \(\pi\) is an isomorphism, so \(\overline A\)
is the automorphism induced by \(\tau\).
This finite action is an image of the pro-\(2\) group
\(\Gamma\simeq\mathbb Z_2\), and consequently has \(2\)-power order.
In characteristic two it is unipotent:
\((\overline A-1)^{2^b}=0\) for some \(b\);
its minimal polynomial divides \((T-1)^r\).
Thus \((\overline A-1)^r=0\).
Since \(r\leq24<32\), we have
\(\overline A^{32}=1\).
Also \(\det\overline A\ne0\), so \(A\in\mathrm{GL}_r(\mathbb Z_2)\).

Write \(A^{32}=1+2B\).
Squaring gives \(A^{64}=1+4(B+B^2)\), and squaring again gives
\begin{equation}\label{eq:sr-A128}
                   A^{128}\equiv1\pmod8.
\end{equation}
Applying \(\pi\) proves the same image congruence for \(\tau\):
\((\tau^{128}-1)M\subseteq8M\).
For \(T_a=5^a\tau\) it follows that
\begin{equation}\label{eq:sr-twist128}
                  (T_a^{128}-1)M\subseteq8M ,
\end{equation}
because \(5^{128a}-1\) is divisible by eight.
The lift is used to prove these congruences; division by two in
an endomorphism of the possibly torsion module \(M\) has not been
assumed.

\subsection{Growth from the fixed layer}
\begin{srproposition}\label{prop:sr-coker-growth}
For every \(a\geq1,m\geq0\),
\[
 c_m(a)\leq c_7(a)+24(m-7)_+ .
\]
If \(c_7(a)\) is finite, every \(c_m(a)\) is finite.
\end{srproposition}
\begin{proof}
For \(m<7\), the factorization of \(T_a^{128}-1\) by
\(T_a^{2^m}-1\) shows that the latter cokernel is a quotient
of the former. Hence \(c_m(a)\leq c_7(a)\).

For \(m=7+k\), put \(V=T_a^{128}\) and
\(S_k=1+V+\cdots+V^{2^k-1}\).
Then \(T_a^{2^m}-1=(V-1)S_k\).
For \(k=0\), \(S_k=1\).
For \(k\geq1\), let \(N=V-1\); by
\eqref{eq:sr-twist128}, \(N^jM\subseteq2^{3j}M\).
The finite binomial identity gives
\[
 S_k=\sum_{j=0}^{2^k-1}{2^k\choose j+1}N^j .
\]
For \(j\geq1\),
\[
 v_2\binom{2^k}{j+1}\geq k-v_2(j+1),\qquad
                       v_2(j+1)\leq j .
\]
The first inequality follows from
\(\binom{2^k}{j+1}=\frac{2^k}{j+1}\binom{2^k-1}{j}\);
the second follows from \(j+1\leq2^j\).
Every nonconstant term consequently maps \(M\) into
\(2^{k+1}M\). Thus
\((S_k-2^k)M\subseteq2^{k+1}M\).

The compact quotient \(Q=M/S_kM\) satisfies
\(2^kQ\subseteq2^{k+1}Q\).
The reverse inclusion is automatic, so iteration and
\(\bigcap_n2^nQ=0\) give \(2^kQ=0\).
Since \(Q\) has at most 24 generators, it is a quotient of
\((\mathbb Z/2^k)^{24}\), and its length is at most \(24k\).

Finally, for any two endomorphisms \(B,C\) of a module, there is an
exact sequence
\[
 0\longrightarrow B(M)/BC(M)
 \longrightarrow M/BC(M)\longrightarrow M/B(M)\longrightarrow0,
\]
and \(B(M)/BC(M)\) is a quotient of \(M/C(M)\).
Apply this with \(B=V-1,C=S_k\).
Length additivity yields \(c_m(a)\leq c_7(a)+24k\),
including in extended lengths, and proves the finiteness assertion.
\end{proof}

\subsection{At most 24 centers govern all the other twists}
Use \(v_2(2)=1\) on finite extensions of \(\mathbb Q_2\).
For a matrix with entries in \(8\mathbb Z_2\), its logarithm and
exponential converge, and the usual inverse identities hold.
Here are the estimates needed to use these functions.
For \(v_2(z)\geq2\),
\[
 v_2(\exp z-1)=v_2(z).
\]
Indeed \((\exp z-1)/z=1+\sum_{n\geq2}z^{n-1}/n!\);
every other term has valuation at least
\(2(n-1)-v_2(n!)\geq n-1\geq1\).
The identical valuation estimates prove convergence for matrices
in the maximum-entry norm.
The formal identities for logarithm and exponential can then be
substituted, because all the relevant power series and their
products converge in this complete matrix algebra.
In particular \(\exp(\log U)=U\) for \(U\equiv1\pmod8\),
and exponentials of commuting matrices multiply.
Also \(v_2(\log5)=2\), since the first term of
\(\log(1+4)\) is four and every subsequent term is divisible by eight.

Take \(U=A^{128}\) and the matrix
\[
              S=-\frac{\log U}{128\log5}\in M_r(\mathbb Q_2).
\]
Choose a finite extension \(K/\mathbb Q_2\) splitting its
characteristic polynomial, and list the eigenvalues
\(s_1,\ldots,s_r\in K\) with multiplicity.
They depend on \(q,d\) and the one chosen lift \(A\), not on the
twist \(a\), the layer \(m\), or a subsequent window.

\begin{srproposition}\label{prop:sr-centers}
If \(B\geq0\) is an integer and
\(|a-s_i|_2\geq2^{-B}\) for every \(i\), then
\begin{equation}\label{eq:sr-center-bound}
                  c_m(a)\leq24(\max(m,7)+2+B)
                         \qquad(m\geq0).
\end{equation}
In particular these cokernels are finite.
\end{srproposition}
\begin{proof}
First let \(m\geq7\).
On the free module, the relevant endomorphism is
\[
 (5^aA)^{2^m}-1
 =\exp\left(2^{m-7}(128a\log5+\log U)\right)-1 .
\]
The matrix in parentheses has entries divisible by eight.
Its eigenvalues, after triangularization over \(K\), are
\(2^m\log5\,(a-s_i)\), and have valuation at least three.
One can check this last assertion without making an isometric
change of basis: a matrix divisible by eight has characteristic
roots divisible by eight, because the characteristic roots of
an integral matrix are integral.
The scalar exponential estimate therefore gives
\[
 v_2\det((5^aA)^{2^m}-1)
       =\sum_{i=1}^r\bigl(m+2+v_2(a-s_i)\bigr)
       \leq r(m+2+B).
\]
The assumed lower distances make every eigenfactor nonzero.
Over the discrete valuation ring \(\mathbb Z_2\), row and column
elimination at an entry of least valuation diagonalizes a matrix
of nonzero determinant.
Its cokernel is finite and has length the valuation of that
determinant: the diagonal factors \(2^{b_i}\) contribute
\(\sum b_i=v_2\det\).
The surjection \(\pi\) induces a surjection from this free
cokernel onto the cokernel defining \(c_m(a)\).
This proves the required bound for \(m\geq7\).
For \(m<7\), use the quotient comparison with \(c_7(a)\)
from the preceding proposition.
\end{proof}

Fix an integer horizon \(Z\geq1\), and put
\(B_Z=\lfloor\log_2 Z\rfloor\).
Define
\begin{equation}\label{eq:sr-exception-set}
 \mathcal E_Z=
 \bigcup_{i=1}^r\{a\in\mathbb Z:1\leq a\leq Z,\
                              |a-s_i|_2<2^{-B_Z}\}.
\end{equation}
Each set in the union has at most one member.
If distinct \(a,b\in[1,Z]\) belonged to it, ultrametricity would give
\(|a-b|_2<2^{-B_Z}\).
But \(0<|a-b|<2^{B_Z+1}\) as ordinary integers, so
\(v_2(a-b)\leq B_Z\), a contradiction.
Thus \(|\mathcal E_Z|\leq r\leq24\), and
\begin{equation}\label{eq:sr-outside-exception}
 c_m(a)\leq24(\max(m,7)+2+B_Z)
       \quad(1\leq a\leq Z,\ a\notin\mathcal E_Z).
\end{equation}
The set is chosen once for the whole horizon.
No new exceptional set is chosen when the real scale varies.

\subsection{The fixed-depth discriminant}
At depth seven,
\[
 L_7=\mathbb Q(\zeta_{512},\sqrt{-q},\sqrt{-d}),\qquad
                         [L_7:\mathbb Q]=1024.
\]
The quadratic radicals are independent over
\(\mathbb Q(\zeta_{512})\) by their odd valuations at \(q,d\).
For later use we prove the explicit bound
\begin{equation}\label{eq:sr-L7-discriminant}
            |\operatorname{disc}L_7|
                     \leq2^{8192}(qd)^{512}.
\end{equation}
The power basis for \(\zeta_{512}\) has discriminant of absolute
value \(2^{2048}\): its polynomial is \(X^{256}+1\), and the
product of the derivatives \(256\zeta^{255}\) at all roots has
absolute value \(256^{256}\).
The quadratic integral bases for the other two fields have
discriminants \(-q,-d\).
Products of these three bases form an integral basis of a full-rank
order in \(L_7\), since the three fields are linearly disjoint.
The trace matrix is the tensor product of the three trace matrices.
Its determinant is
\[
       (2^{2048})^4 q^{512}d^{512}
                          =2^{8192}(qd)^{512}.
\]
For an integral full-rank order of index \(I\) in the maximal
order, the discriminants differ by \(I^2\), by the determinant
of the basis change.
Dropping this index proves \eqref{eq:sr-L7-discriminant}.
In particular, if \(x\geq1\) and \(\log(qd)\leq x^{\epsilon/4}\),
\begin{equation}\label{eq:sr-L7-logdisc}
 \log(2|\operatorname{disc}L_7|)
 \leq (8193\log2+512)x^{\epsilon/4}.
\end{equation}

\subsection{The one high-twist input still required}
Let \(S_2(F)\) be the primes above two in a number field \(F\).
Write
\[
 \mathcal H_2(F,a)=
 H^2(G_{F,S_2(F)},\mathbb Z_2(a+1)).
\]
We use continuous cohomology for the compact coefficient module;
equivalently one uses the inverse finite-coefficient system with
its continuous transition maps.
The classical Poitou--Tate duality background gives, for totally
imaginary \(F\), the pairing
\[
 T_{2,a}(F)\ \times\
 \ker\left(\mathcal H_2(F,a)\longrightarrow
       \prod_{v\mid2}H^2(F_v,\mathbb Z_2(a+1))\right)
                      \longrightarrow\mathbb Q_2/\mathbb Z_2 .
\]
It is a perfect Pontryagin pairing.
This is the same global and local Tate duality input used in
Section~\ref{sec:minus-characterwise}, with mutually dual
coefficients of twists \(-a\) and \(a+1\).
There are no real-place corrections for these fields.
The finite-coefficient pairings pass to the stated limits:
the finite-level groups are finite, hence their inverse systems
satisfy the Mittag--Leffler condition; direct limits are exact,
and only finitely many local places occur.
Consequently finiteness of \(\mathcal H_2(F,a)\) gives
\begin{equation}\label{eq:sr-Sha-H2}
 s_{2,F}(a)\log2\leq\log\#\mathcal H_2(F,a).
\end{equation}
This inequality uses the inclusion of the displayed kernel into
\(\mathcal H_2(F,a)\), not an equality with the whole \(H^2\).

\subsection{The local cohomology terms throughout the tower}
For these particular fields the local correction in global duality can
be computed exactly. This removes a further imported estimate without
proving the remaining global class-counting bound.

\begin{srproposition}\label{prop:sr-H2-exact}
For every \(m\geq0\), the field \(F=L_m\) has exactly two primes above two.
For every
\(a\geq1\), the group \(\mathcal H_2(F,a)\) is finite if and only if
the strict group \(T_{2,a}(F)\) is finite. When they are finite,
\begin{equation}\label{eq:sr-H2-exact}
 \operatorname{length}_{\mathbb Z_2}\mathcal H_2(F,a)
                 =s_{2,F}(a)+m+2+v_2(a).
\end{equation}
\end{srproposition}
\begin{proof}
First determine the completions. Put \(b=m+2\geq2\) and \(C=\mathbb Q(\zeta_{2^b})\).
The polynomial of \(\zeta_{2^b}-1\) is
\((Y+1)^{2^{b-1}}+1\). Its constant term is two, every other
nonleading coefficient is divisible by two, and the leading coefficient
is one. The divisibility of the intermediate binomial coefficients
follows by reducing \((1+Y)^{2^{b-1}}\) modulo two, where it is
\(1+Y^{2^{b-1}}\). Thus the polynomial is Eisenstein at two.
Consequently \(C\) has a single dyadic completion
\(C_2=\mathbb Q_2(\zeta_{2^b})\), totally ramified of degree \(2^{b-1}\).

An odd \(2\)-adic unit congruent to one modulo eight is a square in
\(\mathbb Q_2\). To see sufficiency directly, start with an odd square
root modulo \(2^3\). Given an odd \(b_n\) with
\(b_n^2\equiv u\pmod {2^n}\), \(n\geq3\), choose
\(b_{n+1}=b_n+\delta 2^{n-1}\), \(\delta\in\{0,1\}\).
Modulo \(2^{n+1}\), the square changes by \(\delta b_n2^n\).
Exactly one choice gives the next congruence. These choices are Cauchy
and their limit squares to \(u\).
Necessity follows by squaring an odd integer modulo eight.
Since \(d\equiv7\pmod {16}\), \(-d\equiv1\pmod8\).
Hence adjoining \(\sqrt{-d}\) splits over \(C_2\).

On the other hand \(-q\equiv5\pmod8\).
For \(u=-q\), the polynomial
\[
                 Z^2-Z+(1-u)/4
\]
reduces to \(Z^2+Z+1\) over \(\mathbb F_2\).
It is irreducible and separable. Its root generates the unramified
quadratic extension \(U/\mathbb Q_2\).
For completeness, the ring obtained by adjoining this root to
\(\mathbb Z_2\) is free of rank two and its reduction is
\(\mathbb F_4\). An element whose reduction is nonzero is invertible:
lift its inverse modulo two and correct by the convergent geometric
series. Every nonzero element is a power of two times such a unit,
by factoring the least valuation of its two coefficients.
This ring is therefore a discrete valuation ring with uniformizer two
and residue field \(\mathbb F_4\). It is integrally closed, hence is
the full ring of integers of its fraction field. This proves the
unramified assertion.

A totally ramified extension and an unramified extension have trivial
intersection; after unramified base change \(C_2U/C_2\) has degree
two. Since \(F=C(\sqrt{-q},\sqrt{-d})\), its dyadic algebra is
\[
              F\otimes_{\mathbb Q}\mathbb Q_2
                     \simeq C_2U\times C_2U .
\]
Thus there are exactly two dyadic primes and their completions are
both \(C_2U\), of degree \(2^b\) over \(\mathbb Q_2\).
Here \([F:\mathbb Q]=2^{b+1}\), by the independence of the two
odd radicals already used above.

We next determine the cyclotomic images. Every finite extension in
\(\mathbb Q_2(\mu_{2^\infty})/C_2\) is totally ramified, whereas
\(C_2U/C_2\) is unramified. Their intersection is \(C_2\).
It follows that the cyclotomic image of each local absolute Galois
group is
\[
                    1+2^b\mathbb Z_2 .
\]
The global image is the same. Indeed \(F/C\) is biquadratic.
Each of its three quadratic intermediate fields is ramified at
\(q\) or \(d\), since the corresponding radical
\(-q,-d\), or \(qd\) has odd valuation there and \(C/\mathbb Q\)
is unramified at these odd primes.
No such field lies in the cyclotomic \(2\)-power tower, which is
unramified outside two.
As the cyclotomic tower over \(C\) is procyclic, a nontrivial
intersection with \(F\) would contain a quadratic intermediate
field. Thus \(F\cap\mathbb Q(\mu_{2^\infty})=C\), proving the
global image assertion.

The group \(1+2^b\mathbb Z_2\) is topologically generated by
\(1+2^b\). One direct verification uses logarithm: it identifies
this group with \(2^b\mathbb Z_2\), and
\(\log(1+2^b)/2^b\) is a unit.
For \(a\geq1\),
\[
                 v_2((1+2^b)^a-1)=b+v_2(a).
\]
For odd \(a\) the first binomial term has valuation \(b\) and all
subsequent terms have larger valuation.
Replacing an odd exponent by twice that exponent raises the valuation
by one, using \(z^2-1=(z-1)(z+1)\) with \(z\equiv1\pmod4\).
Iteration proves the formula for every \(a\).
The same valuation holds for the inverse power minus one.

Let \(D_a=\mathbb Q_2/\mathbb Z_2(-a)\).
Its invariants under any of the global or local groups just
considered are therefore cyclic of order
\(2^{b+v_2(a)}\). A class in \(\mathbb Q_2/\mathbb Z_2\)
is fixed precisely when it is killed by the action of the dense
generator minus one; continuity then gives invariance under the
whole group. Local Tate duality identifies
\(H^2(F_v,\mathbb Z_2(a+1))\) with the Pontryagin dual of these
local invariants, so it has length \(b+v_2(a)\).

Finally use the following tail of global Poitou--Tate duality:
\[
 0\longrightarrow T_{2,a}(F)^\vee
 \longrightarrow \mathcal H_2(F,a)
 \longrightarrow \bigoplus_{v\mid2}
                    H^2(F_v,\mathbb Z_2(a+1))
 \longrightarrow H^0(F,D_a)^\vee
 \longrightarrow0 .
\]
Here \({}^\vee\) denotes Pontryagin duality.
The leftmost identification is the strict pairing printed above.
The final map is dual to the diagonal map on invariants, which is
injective. The totally imaginary hypothesis removes archimedean
terms. As for the preceding pairing, this is the compact-coefficient
tail obtained from the finite-coefficient duality sequence.
The other three groups in this sequence are finite.
It follows at once that its first two groups are finite
simultaneously. In that case length additivity gives
\[
 s_{2,F}(a)+2(b+v_2(a))-(b+v_2(a))
                  =s_{2,F}(a)+m+2+v_2(a),
\]
as asserted.
\end{proof}

\begin{srproposition}\label{prop:sr-H2-growth}
If \(\mathcal H_2(L_7,a)\) is finite, then every
\(\mathcal H_2(L_m,a)\) is finite, and
\[
 \operatorname{length}_{\mathbb Z_2}\mathcal H_2(L_m,a)
 \leq
 \operatorname{length}_{\mathbb Z_2}\mathcal H_2(L_7,a)
                                      +25(m-7)_+ .
\]
\end{srproposition}
\begin{proof}
Finiteness at depth seven gives finiteness of \(c_7(a)\) by
Proposition~\ref{prop:sr-H2-exact}. Proposition~\ref{prop:sr-coker-growth}
then gives finiteness of every \(c_m(a)\), and the exact local
calculation gives finiteness of the corresponding \(H^2\).
For \(m\geq7\), add \(m+2+v_2(a)\) to
\(c_m(a)\leq c_7(a)+24(m-7)\), and apply
\eqref{eq:sr-H2-exact} at both depths. This gives \(25(m-7)\).
For \(m<7\), \(c_m(a)\leq c_7(a)\) and \(m+2\leq9\),
giving the asserted bound with zero extra term.
\end{proof}

In particular, the missing size theorem is a global strict-group
estimate, rather than an additional local correction.
More precisely, suppose a uniform \(C\geq1\) gave
\[
 s_{2,L_7}(a)\log2\leq
 \exp\bigl((2048a)^C\log(2|\operatorname{disc}L_7|)\bigr)
\]
together with finiteness, for every pair and every \(a\geq1\).
Then \(K_2(C+1)\) would follow. Set
\(B=(2048a)^C\log(2|\operatorname{disc}L_7|)\).
We have \(B\geq2048a\log2\geq1\), and
\((9+v_2(a))\log2\leq B\), since \(v_2(a)\leq a\).
By \eqref{eq:sr-H2-exact},
\[
 \log\#\mathcal H_2(L_7,a)\leq e^B+B
        \leq e^{2B}
        \leq
 \exp\bigl((2048a)^{C+1}\log(2|\operatorname{disc}L_7|)\bigr).
\]
For the middle inequality, \(B\leq e^B\) and
\(2e^B\leq e^{2B}\) when \(B\geq1\).
Conversely \(K_2(C_K)\) bounds the strict length by
\eqref{eq:sr-Sha-H2}. Neither direction proves the global strict
estimate itself.

\begin{srdefinition}[The specialized high-twist bound]
\label{def:sr-K2}
The assertion \(K_2(C_K)\), for a fixed real \(C_K\geq1\), says
that for every principal pair \(q,d\) and every integer \(a\geq1\),
the group \(\mathcal H_2(L_7,a)\) is finite and
\begin{equation}\label{eq:sr-K2}
 \log\#\mathcal H_2(L_7,a)
 \leq
 \exp\bigl((2048a)^{C_K}
                \log(2|\operatorname{disc}L_7|)\bigr).
\end{equation}
\end{srdefinition}
The archive's general Lemma \(K\) implies this assertion by the
degree calculation above and \(\mu_4\subset L_7\).
Its class-counting and geometric proof has not been replaced by
a standalone proof here. Thus \(K_2(C_K)\) remains an explicit
input to the next theorem.

\subsection{The dyadic estimate on the same common windows}
\begin{srtheorem}\label{thm:sr-dyadic-window}
Assume \(K_2(C_K)\).
Fix \(C_t=401\), \(0<\epsilon<1/52\), and put
\[
 \eta=\frac{\epsilon}{2C_K},\qquad
 J=2^{26}(2C_KC_t/\epsilon)^{24}.
\]
For sufficiently large \(X\), uniformly in the principal pair
\(q,d\) satisfying
\(\log(qd)\leq X^{\epsilon/(4J)}\), there is \(Y\) with
\[
 X^{1/J}\leq Y,\qquad Y^2\leq X
\]
such that for every \(Y\leq x\leq Y^2\), every natural \(h=h(x)\)
with
\[
 x^{2+8\epsilon}\leq2^h\leq2x^{2+8\epsilon},
                  \qquad t=\lfloor x^{401}\rfloor,
\]
all \(s_{2,E_h}(a)\), \(1\leq a<t\), are finite and
\begin{equation}\label{eq:sr-full-dyadic-window}
                \sum_{a=1}^{t-1}s_{2,E_h}(a)
                          \leq\exp(x^{4\epsilon/5}).
\end{equation}
The threshold depends on \(C_K,\epsilon\), not on \(X,q,d\)
or the subsequent choices of \(x,h\).
An admissible depth always exists: take
\(h=\lceil(2+8\epsilon)\log_2x\rceil\).
\end{srtheorem}
\begin{proof}
For the fixed pair choose the one module, lift and list of centers
above. Take \(Z=\lfloor X^{401}\rfloor\) and form
\(\mathcal E_Z\) from \eqref{eq:sr-exception-set}.
Apply the finite-exception interval selection proved in
\eqref{eq:sr-logavoid}--\eqref{eq:sr-J}, with \(s=2\),
\(R=24\) and this \(\eta\).
Its exponent is \(J/2\); therefore its stronger lower bound
\(Y\geq X^{2/J}\) implies the required \(Y\geq X^{1/J}\).
On the resulting interval,
\[
 a\in\mathcal E_Z,\quad a\leq x^{401}
                         \quad\Longrightarrow\quad a\leq x^\eta .
\]
This one exceptional set is used for the entire interval.

For every \(x\) there, \(x\leq X\), so \(t\leq Z\).
Also
\[
 \log(qd)\leq x^{\epsilon/4},\quad
 B_Z\leq401\log_2 X\leq401J\log_2 x,\quad
 h\leq1+(2+8\epsilon)\log_2 x.
\]
For a twist \(a\notin\mathcal E_Z\), combine
\eqref{eq:sr-outside-exception} with
\eqref{eq:sr-dyadic-return}.
Summing over fewer than \(x^{401}\) such twists contributes at most
\[
                     C_1 x^{401}(1+\log x),
\]
where \(C_1\) depends only on \(\epsilon,C_K\) through \(J\).

For an exceptional twist in the range, \(a\leq x^\eta\).
By \eqref{eq:sr-K2} and~\eqref{eq:sr-L7-logdisc},
\[
 \log\#\mathcal H_2(L_7,a)
 \leq\exp(D_0x^{3\epsilon/4}),\qquad
 D_0=2048^{C_K}(8193\log2+512).
\]
Indeed \(a^{C_K}\leq x^{\epsilon/2}\).
Once \(x^{\epsilon/40}\geq D_0\), this is at most
\(\exp(x^{31\epsilon/40})\).
Equations~\eqref{eq:sr-Sha-H2} and
\eqref{eq:sr-dyadic-coker}, followed by
Proposition~\ref{prop:sr-coker-growth} and
\eqref{eq:sr-dyadic-return}, now give
\[
 s_{2,E_h}(a)
 \leq1+\frac{\exp(x^{31\epsilon/40})}{\log2}
                          +24(h-7)_+ .
\]
There are at most 24 exceptional twists.
All these lengths are finite by \(K_2(C_K)\).

The total is therefore at most
\[
 C_1x^{401}(1+\log x)
 +24+\frac{24}{\log2}\exp(x^{31\epsilon/40})
                              +576h .
\]
Dividing this expression by \(\exp(x^{4\epsilon/5})\) gives a
quantity tending to zero:
the logarithm of the polynomial terms is \(O(\log x)\), and
\(31\epsilon/40<4\epsilon/5\).
Thus \eqref{eq:sr-full-dyadic-window} holds for all sufficiently
large \(x\), uniformly in the pair.
Increasing the lower threshold for \(X\) ensures that every
\(x\geq X^{1/J}\) meets all those thresholds.
This proves the theorem and all its quantifiers.
\end{proof}

This proves the dyadic tower growth, finite-center estimate,
24-element exceptional-set construction, discriminant comparison
and common-window deduction.
The high-twist bound \(K_2(C_K)\) is still required.
The uniform plus estimates and the geometric supply assertion
are separate inputs; none follows from the dyadic calculation.

